\section{Uproszczenie wyrażenia logicznego}
\label{sec:uproszczenie}

Wychodzimy od równania~\eqref{eq:main}:
\begin{align}
    Q_{\text{next}}
    &= (P \oplus Q)\land \lnot T \nonumber\\
    &= (P\lnot Q + \lnot P Q)\,\lnot T \nonumber\\
    &= P\,\lnot Q\,\lnot T \;+\; \lnot P\,Q\,\lnot T .
    \label{eq:sop}
\end{align}

Wyrażenie~\eqref{eq:sop} jest \emph{minimalną} postacią
dysjunkcyjną (SOP). Nie da się jej dalej zredukować bez zmiany
interpretacji fizycznej, ponieważ:
\begin{itemize}
    \item pierwszy składnik $P\bar Q\bar T$ odpowiada sytuacji
          ,,światło zgaszone, przycisk wciśnięty, brak przekroczenia czasu''
          – włączenie światła,
    \item drugi składnik $\bar P Q\bar T$ odpowiada sytuacji
          ,,światło świeci, przycisk wciśnięty, brak przekroczenia czasu''
          – wyłączenie światła.
\end{itemize}

Równoważnie można pozostawić postać z bramką XOR:
\begin{equation}
    Q_{\text{next}} = (P \oplus Q)\,\lnot T,
\end{equation}
która jest korzystniejsza sprzętowo (jedna bramka XOR + jedna AND
zamiast czterech bramek AND, dwóch NOT i jednej OR).

\paragraph{Wniosek.} Realizacja sprzętowa powinna użyć
\textbf{bramki XOR} – pozwala to zaoszczędzić 4 bramki względem
postaci SOP.